% Generated draft: requires mathematical and source review.
\documentclass[11pt]{article}
\usepackage[margin=25mm]{geometry}
\usepackage{fontspec}
\setmainfont{Latin Modern Roman}
\usepackage{amsmath,amssymb,mathtools,unicode-math}
\setmathfont{Latin Modern Math}
\usepackage{enumitem,needspace,xurl,hyperref}
\setlength{\parindent}{0pt}
\setlength{\parskip}{5pt}
\setlength{\emergencystretch}{3em}
\newcommand{\N}{\mathbb{N}}
\newcommand{\Z}{\mathbb{Z}}
\newcommand{\Q}{\mathbb{Q}}
\newcommand{\R}{\mathbb{R}}
\newcommand{\C}{\mathbb{C}}
\newcommand{\F}{\mathbb{F}}
\newcommand{\A}{\mathbb{A}}
\newcommand{\PP}{\mathbb{P}}
\newcommand{\E}{\mathbb{E}}
\newcommand{\eps}{\varepsilon}
\newcommand{\dd}{\,\mathrm{d}}
\begin{document}
\section*{UnsolvedMath ID: 801169}
\textbf{Source problem code:} catalogue entry
\subsection{Definitions and mathematical statement}
For \(0<b<1\), let
\[
E_b=(-1,-b)\cup(b,1),\qquad \widehat{\psi_b}=\mathbf{1}_{E_b},
\]
where \(\mathbf{1}_{E_b}\) denotes the indicator function of \(E_b\), and \(\psi_b\in L^2(\R)\).

The supplied excerpt does not specify the Fourier transform normalization or the dilation and translation conventions for a frame wavelet. Consequently, the following explicit formulation is conditional on the standard dyadic, integer-translation interpretation with Fourier transform
\[
\widehat f(\xi)=\int_{\R}f(x)e^{-2\pi i x\xi}\,\dd x
\]
for \(f\in L^1(\R)\cap L^2(\R)\), extended to \(L^2(\R)\). Under this interpretation, define
\[
\psi_{b,j,k}(x)=2^{j/2}\psi_b(2^j x-k),\qquad j,k\in\Z.
\]
The function \(\psi_b\) is a frame wavelet if there exist real constants \(0<A_b\le B_b<\infty\) such that, for every \(f\in L^2(\R)\),
\[
A_b\|f\|_2^2
\le \sum_{j\in\Z}\sum_{k\in\Z}
\left|\langle f,\psi_{b,j,k}\rangle\right|^2
\le B_b\|f\|_2^2,
\qquad
\langle f,g\rangle=\int_{\R}f(x)\overline{g(x)}\,\dd x.
\]
Determine exactly the set
\[
\left\{b\in\left(\frac18,\frac16\right]:
\psi_b\text{ is a frame wavelet}\right\}.
\]
The lower endpoint \(1/8\) is excluded and the upper endpoint \(1/6\) is included. The Fourier and affine-system conventions above require confirmation against the original source before this conditional formulation can be accepted as the intended statement.
\subsection{Short English statement}
Determine precisely which parameters \(1/8<b\le 1/6\) make the function with Fourier transform \(\mathbf{1}_{(-1,-b)\cup(b,1)}\) a frame wavelet for \(L^2(\R)\). The explicit frame inequalities use a stated conventional interpretation whose agreement with the original source remains to be checked.
\subsection{Sources}
\begin{itemize}
\item Supplied local UnsolvedMath registry record: permanent ID \(801169\), source problem code \texttt{catalogue entry}, titled \emph{Open Problems in Wavelet Theory --- Problem 8}.
\item Bownik and Rzeszotnik, \emph{Open Problems in Wavelet Theory} (2019), Problem 8. Supplied source URL: \url{https://arxiv.org/abs/1905.08230}. The record reports access on 2026-10-16 and extraction from arXiv LaTeX; these are supplied provenance claims, not independently verified access.
\item The supplied problem text attributes the problem to Weber and the first author, Bownik, with reference label [BW]. No bibliographic details for [BW] are supplied.
\item The registry labels the problem \emph{open}. This is a supplied status label, not an independently verified assessment of the literature.
\end{itemize}
\end{document}
